% =====================================================================
%  3.1 Funksjonsbegrepet og lineære funksjoner — fasit
% =====================================================================

\fk{Del 1 — Hva er en funksjon?}
\fo{1.1} a) $-5$ \ b) $7$ \ c) $-11$
\fo{1.2} a) $3$ \ b) $x = -1$ og $x = 1$ \ c) $V_f = \{-1, 0, 3\}$
\fo{1.3} a) $x \neq -5$, altså $\langle\leftarrow, -5\rangle \cup \langle -5, \rightarrow\rangle$ \ b) $[1, \rightarrow\rangle$ \ c) alle tall
\fo{1.4} a) $3$ \ b) $f(2a) = 8a^2 - 2a$ \ c) $f(a+2) = 2a^2 + 7a + 6$
d) Nei. $2f(a) = 4a^2 - 2a$, og $f(2a) - 2f(a) = 4a^2$. De er like bare når $a = 0$.
\fo{1.5} a) $6 - 2x \ge 0$, så $\langle\leftarrow, 3]$
b) $x^2 - 9 \neq 0$, så $x \neq \pm 3$: $\langle\leftarrow, -3\rangle \cup \langle -3, 3\rangle \cup \langle 3, \rightarrow\rangle$
c) Nevneren må være positiv (rot, og ikke 0): $x + 1 > 0$, så $\langle -1, \rightarrow\rangle$
\fo{1.6} Feilen er i overgangen til $a^2 - 2a + 1 - 4a - 4$: $-4(a-1) = -4a + 4$, ikke $-4a - 4$. Riktig: $a^2 - 2a + 1 - 4a + 4 = a^2 - 6a + 5$.
\fo{1.7} a) $h(4) = -3$ \ b) $5 - 2x = 9$ gir $x = -2$
c) I a) er $x$ gitt, og vi finner funksjonsverdien. I b) er funksjonsverdien gitt, og vi finner $x$. a) er en utregning, b) er en likning.
d) $2x+1$: uttrykk. $2x+1=7$: likning (sann bare for $x=3$). $2(x+1)=2x+2$: identitet (sann for alle $x$). $f(x)=2x+1$: funksjon (en regel som gir én verdi for hver $x$).
\fo{1.8} a) Ja. Hver elev har nøyaktig én fødselsdato.
b) Nei. Flere elever kan ha samme fødselsdato, så én dato kan gi flere elever.
c) Ja. Hver $x$ gir nøyaktig ett kvadrat.
d) Nei. For $x = 4$ er både $y = 2$ og $y = -2$ riktig.
\fo{1.9} a) Ett eksempel: $f(x) = 2x$ med $D_f = [1, 5]$. Da er $f(1) = 2$ og $f(5) = 10$, og alle verdier mellom blir tatt. Flere svar er riktige.
b) Ja, uendelig mange. For eksempel $g(x) = \frac{1}{2}(x-1)^2 + 2$ med $D_g = [1, 5]$: $g(1) = 2$, $g(5) = 10$, og $g$ stiger i hele intervallet.
c) Ja, det er $h(x) = -2x + 12$: $h(1) = 10$ og $h(5) = 2$, og alle verdier mellom 2 og 10 blir tatt.
\fo{1.10} Rota krever $x - a \ge 0$. Intervallet starter i 3, så $a = 3$. Nevneren er 0 for $x = b$, og 7 er utelatt, så $b = 7$.
\fo{Sjekk 1} a) $7$ \ b) $[6, \rightarrow\rangle$ \ c) For $x = 0$ gir sirkelen to $y$-verdier ($y = 1$ og $y = -1$). En funksjon gir bare én verdi per $x$.

\fk{Del 2 — Lineære funksjoner}
\fo{2.1} a) $3$ \ b) $-1$ \ c) $0$
\fo{2.2} a) $a = 4$, $b = -7$ \ b) $a = -\frac{1}{2}$, $b = 3$ \ c) $a = 5$, $b = 0$ \ d) $a = -3$, $b = 6$.
Bare $h$ er proporsjonal ($b = 0$).
\fo{2.3} a) $a = \frac{11 - 5}{4 - 2} = 3$ \ b) $5 = 3 \cdot 2 + b$, så $b = -1$ \ c) $f(x) = 3x - 1$
\fo{2.4} a) $a = \frac{-3 - 9}{1 - (-2)} = -4$, $b = 1$: $f(x) = -4x + 1$
b) $a = \frac{3}{2}$, $b = -1$: $f(x) = \frac{3}{2}x - 1$
c) $a = -\frac{1}{2}$, $b = 3$: $f(x) = -\frac{1}{2}x + 3$
\fo{2.5} a) $a = \frac{180 - 216}{10 - 4} = -6$ og $b = 240$: $f(x) = -6x + 240$
b) $a = -6$: tanken mister 6 liter per minutt. $b = 240$: det var 240 liter i tanken da den begynte å tømmes ($x = 0$ gir mening her).
c) Etter 25 minutter. d) $f(x) = 0$ gir $x = 40$. $D_f = [0, 40]$ og $V_f = [0, 240]$.
\fo{2.6} Feilen er i overgangen til $8 = -6 + b$: $-3 \cdot (-2) = +6$. Riktig: $8 = 6 + b$, så $b = 2$ og $f(x) = -3x + 2$. Kontroll: $f(2) = -4$.
\fo{2.7} a) Noen ganger. $f(s+t) - f(s) - f(t) = a(s+t) + b - as - b - at - b = -b$. Det stemmer bare når $b = 0$.
b) Noen ganger. $f(2x) - 2f(x) = -b$, så bare når $b = 0$.
c) Noen ganger. Hvis $a \neq 0$: ett punkt, $x = -\frac{b}{a}$. Hvis $a = 0$ og $b \neq 0$: aldri. Hvis $a = 0$ og $b = 0$: grafen er hele $x$-aksen.
\fo{2.8} a) $a = 60\,000$: prisen øker med 60 000 kr per m$^2$.
b) $P(0) = -400\,000$. En leilighet på 0 m$^2$ finnes ikke, og prisen er negativ. Her er $b$ bare et regnetall, ikke en startverdi.
c) $P(30) = 1\,400\,000$ og $P(120) = 6\,800\,000$. Funksjonen stiger, så $V_P = [1\,400\,000, 6\,800\,000]$.
d) Nei. $P(60) = 3\,200\,000$, men $2 \cdot P(30) = 2\,800\,000$. Det måtte vært $b = 0$ (proporsjonal): da er $P(60) = 2P(30)$.
\fo{2.9} a) $a = \frac{5 - (-1)}{3 - 0} = 2$ og $b = -1$. Da er $k = 2 \cdot 7 - 1 = 13$.
b) Linjen er $y = 2x - 1$, og i $x = 7$ er den på $y = 13$. $C$ ligger over linjen når $k > 13$ (under når $k < 13$).
c) På en rett linje er stigningstallet det samme mellom alle punkt (like steg i $x$ gir like steg i $y$).
\fo{Sjekk 2} a) $2$ \ b) $7$ \ c) $f(0) = a \cdot 0 = 0$, så punktet $(0, 0)$ ligger på grafen.

\fk{Ekstra utfordring}
\fo{U1} a) $f(x) = 20x$ for $0 \le x \le 2$. For $x > 2$: $40 + 35(x - 2) = 35x - 30$.
b) Grafen er to rette stykker. Den får en knekk i $x = 2$, der stigningen går fra 20 til 35.
c) $f(5) = 40 + 3 \cdot 35 = 145$ kr.
\fo{U2} a) $N(x) = 19x + 65$ og $B(x) = 14x + 100$
b) 4 km: $N = 141$ kr, $B = 156$ kr. 10 km: $N = 255$ kr, $B = 240$ kr.
c) $N(x) = B(x)$ gir $5x = 35$, $x = 7$. Under 7 km er Nordtaxi billigst, så påstanden stemmer ikke. Bytaxi er billigst først etter 7 km.
\fo{U3} a) \lstinline[style=laereverk,numbers=none,frame=none]|b = y[0] - a * x[0]|
b) Første test: alle stigningstall er 3, og $b = 2 - 3 \cdot 0 = 2$. Utskrift: «Ja: a = 3 og b = 2». Andre test: de to første stigningstallene er 1 og 2, så den skriver «Nei».
c) Desimaltall lagres ikke eksakt. For eksempel er $0{,}3 - 0{,}2$ lik $0{,}09999999999999998$ og ikke $0{,}1$. Da kan to like stigningstall bli ulike. Brøker er eksakte, så vi får $a$ og $b$ som eksakte tall, for eksempel $\frac{1}{3}$ i stedet for $0{,}333\ldots$.
\fo{U4} a) Mellom $(1, 2)$ og $(3, 8)$: $a = \frac{6}{2} = 3$. Mellom $(1, 2)$ og $(-2, -7)$: $a = \frac{-9}{-3} = 3$. Samme stigningstall fra samme punkt, så de ligger på samme linje.
b) $f(x) = 3x - 1$. Kontroll: $f(-2) = -7$.
\fo{U5} a) $\Delta x = 2 - 2 = 0$, og vi kan ikke dele på 0.
b) Linjen er $x = 2$. Den loddrette linjen $x = 2$ treffer grafen i uendelig mange punkt, så $x = 2$ har flere funksjonsverdier.
\fo{U6} a) Hvis $A$, $B$ og $C$ ligger på linjen $y = ax + b$, gir begge stigningstallene $a$, så de er like. Omvendt: hvis stigningstallet mellom $A$ og $B$ og mellom $B$ og $C$ begge er $a$, ligger $C$ på linjen gjennom $B$ med stigning $a$. Den linjen er den samme som linjen gjennom $A$ og $B$.
b) Stigningstallene må være like: $\frac{k-2}{2} = \frac{8-k}{k-3}$ ($k \neq 3$). Da er $(k-2)(k-3) = 2(8-k)$, så $k^2 - 5k + 6 = 16 - 2k$, altså $k^2 - 3k - 10 = 0$ og $k = 5$ eller $k = -2$.
Kontroll: $k = 5$ gir $(1,2)$, $(3,5)$, $(5,8)$ med $a = \frac{3}{2}$. $k = -2$ gir $(1,2)$, $(3,-2)$, $(-2,8)$ med $a = -2$. Punktene er ulike i begge tilfeller.
\fo{U7} Flere svar er riktige.
a) $f(x) = \sqrt{x^2 - 7x + 10} = \sqrt{(x-2)(x-5)}$. Produktet er ikke negativt for $x \le 2$ og $x \ge 5$. En stykkevis funksjon går også, for eksempel $f(x) = x$ for $x \le 2$ og $f(x) = 2x$ for $x \ge 5$.
b) $f(x) = \frac{1}{(x+1)(x-6)}$ ($x \neq -1$ og $x \neq 6$).
c) $f(x) = x^2$, $f(x) = |x|$ eller $f(x) = \sqrt{x}$ har alle $V_f = [0, \rightarrow\rangle$.
